\section{Checklist de examen}

\begin{trampa}{Si cambian los numeros}
Si cambian \(100\), \(1\), \(c\), o se agrega un costo al minorista, no cambia la logica. Cambian los parametros:
\[
p=a-bq,
\qquad
c=\text{costo proveedor},
\qquad
k=\text{costo unitario adicional del minorista}.
\]

Si existe \(k\), reemplaza \(c\) por \(c+k\) en la cadena centralizada y usa:
\[
w_D=\frac{a+c-k}{2},
\qquad
q_D=\frac{a-c-k}{4b},
\qquad
p_D=\frac{3a+c+k}{4}.
\]
Centralizado:
\[
q_C=\frac{a-c-k}{2b},
\qquad
p_C=\frac{a+c+k}{2}.
\]
\end{trampa}

\begin{materia}{Frases que suman puntos}
\begin{itemize}[itemsep=2pt]
    \item ``La cadena descentralizada no maximiza la utilidad total, porque cada agente optimiza localmente.''
    \item ``La doble marginalizacion genera un precio final mayor y una cantidad vendida menor.''
    \item ``El contrato coordina si induce al minorista a elegir la cantidad centralizada.''
    \item ``El parametro \(\alpha\) no cambia la utilidad total si el contrato coordina; solo reparte la utilidad entre agentes.''
    \item ``La eleccion entre contratos depende de eficiencia total y distribucion de beneficios.''
\end{itemize}
\end{materia}

\begin{materia}{Mini algoritmo}
\[
\begin{array}{c|l}
\text{Paso} & \text{Que hacer} \\
\hline
1 & \text{Escribir demanda inversa } p=a-bq.\\
2 & \text{Resolver respuesta del minorista } q(w).\\
3 & \text{Maximizar utilidad del proveedor para obtener } w_D.\\
4 & \text{Calcular } q_D,p_D,\Pi_s^D,\Pi_r^D,\Pi_{SC}^D.\\
5 & \text{Resolver centralizado } q_C,p_C,\Pi_{SC}^C.\\
6 & \text{Comparar: cantidad, precio, utilidad total.}\\
7 & \text{Si hay contrato, verificar si logra } q_C.\\
8 & \text{Repartir utilidad y chequear participacion de ambos.}
\end{array}
\]
\end{materia}

\begin{materia}{Como leer la notacion de utilidad}
La regla es:
\[
\Pi_{\text{quien}}^{\text{escenario}}.
\]
El subindice dice quien gana la utilidad:
\[
s=\text{proveedor},
\qquad
r=\text{minorista},
\qquad
SC=\text{cadena completa}.
\]
El superindice dice en que escenario:
\[
D=\text{descentralizado},
\qquad
C=\text{centralizado},
\qquad
A,B=\text{contratos comparados}.
\]
Ejemplos:
\[
\Pi_s^D=\text{utilidad del proveedor sin coordinacion},
\]
\[
\Pi_r^D=\text{utilidad del minorista sin coordinacion},
\]
\[
\Pi_{SC}^C=\text{utilidad maxima total si la cadena se coordina perfectamente}.
\]
\end{materia}

\begin{trampa}{Errores frecuentes}
\begin{itemize}[itemsep=2pt]
    \item Confundir \(w\) con \(p\). \(w\) es precio mayorista; \(p\) es precio al cliente final.
    \item Maximizar primero al proveedor sin considerar la respuesta del minorista.
    \item Decir que un contrato es mejor solo porque mejora al proveedor. Debe evaluarse tambien al minorista y a la utilidad total.
    \item Olvidar que una tarifa fija \(F\) no cambia la cantidad si no altera el costo marginal del minorista.
    \item Olvidar la frase final de interpretacion: doble marginalizacion.
\end{itemize}
\end{trampa}
